\documentclass[10pt,a4paper]{amsart}
\usepackage{amsmath,amssymb,amsthm,graphicx}
\usepackage{verbatim}
\newenvironment{reference}{\comment}{\endcomment}
%\newenvironment{reference}{}{}
\newenvironment{slogan}{\comment}{\endcomment}
\newenvironment{history}{\comment}{\endcomment}
% For commutative diagrams we use Xy-pic
\usepackage[all]{xy}

% We use 2cell for 2-commutative diagrams.
\xyoption{2cell}
% Theorem environments.
%
\theoremstyle{plain}
\newtheorem{theorem}[subsection]{Theorem}
\newtheorem{proposition}[subsection]{Proposition}
\newtheorem{lemma}[subsection]{Lemma}

\theoremstyle{definition}
\newtheorem{definition}[subsection]{Definition}
\newtheorem{example}[subsection]{Example}
\newtheorem{exercise}[subsection]{Exercise}
\newtheorem{situation}[subsection]{Situation}

\theoremstyle{remark}
\newtheorem{remark}[subsection]{Remark}
\newtheorem{remarks}[subsection]{Remarks}

\numberwithin{equation}{subsection}

% Macros
%
\def\lim{\mathop{\mathrm{lim}}\nolimits}
\def\colim{\mathop{\mathrm{colim}}\nolimits}
\def\Spec{\mathop{\mathrm{Spec}}}
\def\Hom{\mathop{\mathrm{Hom}}\nolimits}
\def\Ext{\mathop{\mathrm{Ext}}\nolimits}
\def\SheafHom{\mathop{\mathcal{H}\!\mathit{om}}\nolimits}
\def\SheafExt{\mathop{\mathcal{E}\!\mathit{xt}}\nolimits}
\def\Sch{\mathit{Sch}}
\def\Mor{\mathop{\mathrm{Mor}}\nolimits}
\def\Ob{\mathop{\mathrm{Ob}}\nolimits}
\def\Sh{\mathop{\mathit{Sh}}\nolimits}
\def\NL{\mathop{N\!L}\nolimits}
\def\CH{\mathop{\mathrm{CH}}\nolimits}
\def\proetale{{pro\text{-}\acute{e}tale}}
\def\etale{{\acute{e}tale}}
\def\QCoh{\mathit{QCoh}}
\def\Ker{\mathop{\mathrm{Ker}}}
\def\Im{\mathop{\mathrm{Im}}}
\def\Coker{\mathop{\mathrm{Coker}}}
\def\Coim{\mathop{\mathrm{Coim}}}

% Boxtimes
%
\DeclareMathSymbol{\boxtimes}{\mathbin}{AMSa}{"02}

%
% Macros for moduli stacks/spaces
%
\def\QCohstack{\mathcal{QC}\!\mathit{oh}}
\def\Cohstack{\mathcal{C}\!\mathit{oh}}
\def\Spacesstack{\mathcal{S}\!\mathit{paces}}
\def\Quotfunctor{\mathrm{Quot}}
\def\Hilbfunctor{\mathrm{Hilb}}
\def\Curvesstack{\mathcal{C}\!\mathit{urves}}
\def\Polarizedstack{\mathcal{P}\!\mathit{olarized}}
\def\Complexesstack{\mathcal{C}\!\mathit{omplexes}}
% \Pic is the operator that assigns to X its picard group, usage \Pic(X)
% \Picardstack_{X/B} denotes the Picard stack of X over B
% \Picardfunctor_{X/B} denotes the Picard functor of X over B
\def\Pic{\mathop{\mathrm{Pic}}\nolimits}
\def\Picardstack{\mathcal{P}\!\mathit{ic}}
\def\Picardfunctor{\mathrm{Pic}}
\def\Deformationcategory{\mathcal{D}\!\mathit{ef}}
\usepackage{cleveref}
\crefname{originallemma}{original lemma}{original lemmas}
\Crefname{originallemma}{Original lemma}{Original lemmas}
\makeatletter
\newlabel{lemma-bijective-map}{{lemma-bijective-map}{source}{}{sourcepointer-1}{}}
\newlabel{lemma-bijective-map@cref}{{[originallemma][0][]lemma-bijective-map}{[1][0][]source}{}{}{}}
\makeatother
\crefname{originallemma}{original lemma}{original lemmas}
\Crefname{originallemma}{Original lemma}{Original lemmas}
\makeatletter
\newlabel{lemma-closed-map}{{lemma-closed-map}{source}{}{sourcepointer-2}{}}
\newlabel{lemma-closed-map@cref}{{[originallemma][0][]lemma-closed-map}{[1][0][]source}{}{}{}}
\makeatother
\crefname{originallemma}{original lemma}{original lemmas}
\Crefname{originallemma}{Original lemma}{Original lemmas}
\makeatletter
\newlabel{lemma-image-open-technical}{{lemma-image-open-technical}{source}{}{sourcepointer-3}{}}
\newlabel{lemma-image-open-technical@cref}{{[originallemma][0][]lemma-image-open-technical}{[1][0][]source}{}{}{}}
\makeatother
\crefname{originallemma}{original lemma}{original lemmas}
\Crefname{originallemma}{Original lemma}{Original lemmas}
\makeatletter
\newlabel{lemma-intersection-empty}{{lemma-intersection-empty}{source}{}{sourcepointer-4}{}}
\newlabel{lemma-intersection-empty@cref}{{[originallemma][0][]lemma-intersection-empty}{[1][0][]source}{}{}{}}
\makeatother
\title{A surjection of quasi-compact Hausdorff spaces over an extremally disconnected target which is minimal is a homeomorphism}
\author{The Stacks Project Authors}
\date{stacks-project:topology.tex:local-snapshot-20261012; GNU FDL 1.2 or later}
\begin{document}
\maketitle
\noindent\textit{Editorial note.} All delivered mathematical text is the author's own, copied from the pinned Stacks Project chapter source sources/stacks-topology/pinned-original-source.tex; no statement or proof step was rewritten, reordered or removed.\par\medskip
\noindent\textit{Editorial note.} Preamble provenance: the Stacks Project keeps one single shared preamble for every chapter; this card takes it directly from the staged sources/stacks-topology/preamble.tex as three verbatim spans (lines 8-12, 14-18 and 154-217), so no derived file was written and nothing was added or redefined. The wrapper class and the base packages are supplied by the assembler.\par\medskip
\noindent\textit{Editorial note.} Bibliography provenance: neither the counted statement nor its proof contains a citation (this lane rescans the whole span before writing the card), so the card carries no bibliography and no entry was reconstructed.\par\medskip
\noindent\textit{Editorial note.} Selection provenance: this card was opened under the file-based control contract (CONTROL.md), with the mechanical pre-card gate run on its statement and proof interval before the card was written.\par\medskip
\noindent\textit{Allowed standard prerequisites (editorial).}\begin{itemize}
\item The statement that a bijective continuous map from a quasi-compact to a Hausdorff space is a homeomorphism. Invoked at the start of the counted proof (source line 5488). Reduces the assertion to injectivity.
\item The statement that a continuous map of quasi-compact spaces with extremally disconnected target is closed. Invoked inside the counted proof (source lines 5492-5493). Makes the images of the two complements closed.
\item The technical containment for the image of an open subset. Invoked inside the counted proof (source lines 5496-5498). Places the common image point in the closures of the two complements.
\item The statement that closures of disjoint opens in an extremally disconnected space are disjoint. Invoked at the end of the counted proof (source lines 5498-5499). Contradicts the two closure memberships.
\end{itemize}\medskip
\ifdefined\hypertarget\hypertarget{sourcepointer-1}{}\fi
\noindent\textit{Source reference.} Original lemma lemma-bijective-map: sources/stacks-topology/pinned-original-source.tex, source lines 3498--3498. Navigation pointer only: the counted text refers to the Stacks item lemma-bijective-map, which is an existing labeled result of this same chapter and is not part of this card. The pointer carries the target's own label as its number and its definition line as its span; its content is not reproduced here.\par\medskip
\ifdefined\hypertarget\hypertarget{sourcepointer-2}{}\fi
\noindent\textit{Source reference.} Original lemma lemma-closed-map: sources/stacks-topology/pinned-original-source.tex, source lines 3474--3474. Navigation pointer only: the counted text refers to the Stacks item lemma-closed-map, which is an existing labeled result of this same chapter and is not part of this card. The pointer carries the target's own label as its number and its definition line as its span; its content is not reproduced here.\par\medskip
\ifdefined\hypertarget\hypertarget{sourcepointer-3}{}\fi
\noindent\textit{Source reference.} Original lemma lemma-image-open-technical: sources/stacks-topology/pinned-original-source.tex, source lines 5449--5449. Navigation pointer only: the counted text refers to the Stacks item lemma-image-open-technical, which is an existing labeled result of this same chapter and is not part of this card. The pointer carries the target's own label as its number and its definition line as its span; its content is not reproduced here.\par\medskip
\ifdefined\hypertarget\hypertarget{sourcepointer-4}{}\fi
\noindent\textit{Source reference.} Original lemma lemma-intersection-empty: sources/stacks-topology/pinned-original-source.tex, source lines 5466--5466. Navigation pointer only: the counted text refers to the Stacks item lemma-intersection-empty, which is an existing labeled result of this same chapter and is not part of this card. The pointer carries the target's own label as its number and its definition line as its span; its content is not reproduced here.\par\medskip
\begin{lemma}
\label{lemma-isomorphism}
Let $f : X \to Y$ be a continuous map of Hausdorff quasi-compact
topological spaces. If $Y$ is extremally disconnected, $f$ is surjective,
and $f(Z) \not = Y$ for every proper closed subset $Z$ of $X$, then
$f$ is a homeomorphism.
\end{lemma}

\begin{proof}
By Lemma \ref{lemma-bijective-map} it suffices to show that $f$ is injective.
Suppose that $x, x' \in X$ are distinct points with $y = f(x) = f(x')$.
Choose disjoint open neighbourhoods $U, U' \subset X$ of $x, x'$.
Observe that $f$ is closed (Lemma \ref{lemma-closed-map}) hence
$T = f(X \setminus U)$ and $T' = f(X \setminus U')$ are closed in $Y$.
Since $X$ is the union of $X \setminus U$ and $X \setminus U'$ we see that
$Y = T \cup T'$. By Lemma \ref{lemma-image-open-technical} we
see that $y$ is contained in the closure of $Y \setminus T$ and
the closure of $Y \setminus T'$. On the other hand, by
Lemma \ref{lemma-intersection-empty},
this intersection is empty. In this way we obtain the desired contradiction.
\end{proof}

\end{document}
